\section{Simulación de 5 cuerpos gravitacionales}

Se simula el movimiento de una partícula $C_5$ bajo la atracción gravitacional
de cuatro masas fijas ubicadas en los vértices de un cuadrado de lado $a = 2$.
El objetivo es encontrar las trayectorias más cercanas a ser órbitas cerradas,
variando la velocidad horizontal inicial de $C_5$.

\subsection{Modelo y aceleraciones}

El sistema consiste en cuatro masas fijas ubicadas en los vértices de un
cuadrado de lado $a = 2$, con coordenadas:

\begin{equation}
    (x_k, y_k) \in \left\{
        \left(\tfrac{a}{2},\tfrac{a}{2}\right),\;
        \left(-\tfrac{a}{2},\tfrac{a}{2}\right),\;
        \left(-\tfrac{a}{2},-\tfrac{a}{2}\right),\;
        \left(\tfrac{a}{2},-\tfrac{a}{2}\right)
    \right\}
\end{equation}

La partícula $C_5$ se desplaza libremente bajo la atracción gravitacional
de las cuatro masas, que permanecen estáticas. La aceleración de $C_5$
en la posición $(x, y)$ se obtiene sumando la contribución de cada masa:

\begin{equation}
    \vec{a}_5 = \sum_{k=1}^{4} \frac{\vec{r}_k - \vec{r}_5}{|\vec{r}_k - \vec{r}_5|^3}
\end{equation}

donde $\vec{r}_k$ es la posición del vértice $k$ y $\vec{r}_5 = (x, y)$ es la
posición actual de $C_5$. Esta expresión corresponde a la ley de gravitación
de Newton con $G = 1$ y masas unitarias.

Se detecta una colisión cuando $C_5$ se aproxima a menos de $R_\text{col} = 0.5$
de alguna masa fija, descartando esa trayectoria. Para explorar distintas
trayectorias se realiza un barrido sobre la velocidad horizontal inicial
$v_{x0} \in [0,\, 1.4]$ con paso $0.02$, manteniendo fijas la posición inicial
$(x_0, y_0) = (2, 7)$ y la velocidad vertical $v_{y0} = 0$. De todas las
trayectorias válidas se seleccionan las cinco cuya trayectoria más se aproxima
al punto de partida, es decir, las más cercanas a ser órbitas cerradas.

\subsection{Implementación}

El primer bloque define los parámetros del sistema: el lado del cuadrado
$a = 2$, el paso de integración $h = 0.01$, el tiempo de simulación
$t_\text{fin} = 750$, las condiciones iniciales de $C_5$ en $(2, 7)$ con
$v_{y0} = 0$, y el barrido de velocidades $v_{x0} \in [0,\, 1.4]$.
También se establecen las posiciones fijas de los cuatro vértices mediante
arreglos \texttt{numpy} y el radio de colisión $R_\text{col} = 0.5$.

\inputminted[linenos, frame=lines, fontsize=\footnotesize,
    breaklines=true, firstline=4, lastline=32
]{python}{code/problem03.py}
\captionof{listing}{Parámetros del sistema, condiciones iniciales y posiciones de los vértices.}
\label{lst:prob3_1}

\vspace{10pt}

El segundo bloque define la función \texttt{acel\_cuadrado}, que recibe la
posición $(x, y)$ de $C_5$ y calcula su aceleración gravitacional sumando la
contribución de las cuatro masas fijas. Para cada masa $k$ se calculan las
diferencias $dx = x - x_k$ y $dy = y - y_k$, la distancia al cubo
$r^3 = (dx^2 + dy^2)^{1.5}$, y se acumula la contribución en $ax$ y $ay$.
Luego, el bucle principal itera sobre cada valor de $v_{x0}$, integra
la trayectoria con el método de Euler, verifica colisiones con cada vértice
y almacena las trayectorias válidas junto con su distancia mínima al punto
de partida. Al finalizar, se ordenan por distancia mínima y se seleccionan
las cinco más cerradas.

\inputminted[linenos, frame=lines, fontsize=\footnotesize,
    breaklines=true, firstline=34, lastline=101
]{python}{code/problem03.py}
\captionof{listing}{Función de aceleración gravitacional, integración por Euler y selección de trayectorias.}
\label{lst:prob3_2}

\vspace{10pt}

El tercer bloque genera la figura con fondo negro. Cada una de las cinco
trayectorias seleccionadas se traza con un color distinto (rojo, verde, azul,
magenta y amarillo) e incluye su valor de $v_{x0}$ en la leyenda. Las cuatro
masas fijas se representan como círculos rellenos en los vértices del cuadrado,
unidos por un contorno punteado blanco, y $C_5$ aparece como un círculo blanco
en el origen como referencia visual.

\inputminted[linenos, frame=lines, fontsize=\footnotesize,
    breaklines=true, firstline=103, lastline=163
]{python}{code/problem03.py}
\captionof{listing}{Configuración visual y generación de la figura de trayectorias.}
\label{lst:prob3_3}

\vspace{10pt}

La Figura \ref{fig:prob3_1} muestra las cinco trayectorias más cerradas
obtenidas para $C_5$ durante el tiempo de simulación $t_\text{fin} = 750$,
correspondientes a $v_{x0} \in \{0.56,\, 0.58,\, 0.60,\, 0.62,\, 0.64\}$.

\begin{figure}[H]
  \centering
  \includegraphics[width=1.0\textwidth]{images/problem03.png}
  \caption{Trayectorias de $C_5$ bajo la atracción de cuatro masas fijas en los
  vértices del cuadrado. Se muestran las cinco órbitas cuya trayectoria más se
  aproxima al punto de partida $(2, 7)$.}
  \label{fig:prob3_1}
\end{figure}

\vspace{10pt}

Se observa que $C_5$ describe trayectorias complejas bajo la influencia
simultánea de las cuatro masas fijas. Ninguna de las órbitas obtenidas es
perfectamente cerrada, pero las seleccionadas presentan el mayor grado de
recurrencia al punto de partida. La simetría del cuadrado impone cierta
regularidad en las trayectorias, que tienden a rodear el conjunto de masas
o a oscilar entre ellas. La velocidad inicial $v_{x0}$ determina si $C_5$
escapa hacia el exterior, queda atrapada en una órbita acotada o impacta
contra alguna de las masas fijas.

\subsection{Conclusión}

La simulación reproduce la dinámica gravitacional de una partícula $C_5$ bajo
la atracción de cuatro masas fijas dispuestas en los vértices de un cuadrado,
mediante la ley de gravitación de Newton con $G = 1$ y masas unitarias. El
método de Euler con paso $h = 0.01$ permite integrar el sistema de forma
sencilla, aunque acumula error con el tiempo, lo que limita la precisión de
las trayectorias para simulaciones largas como $t_\text{fin} = 750$. Las
trayectorias obtenidas muestran que la geometría del cuadrado genera un campo
gravitacional con simetría cuadrada que atrapa a $C_5$ en órbitas complejas
para valores de $v_{x0}$ en el rango $[0.56,\, 0.64]$.